\documentclass[10pt,a4paper]{amsart}
\usepackage[margin=19mm]{geometry}
\usepackage{amsmath,amssymb,mathtools}
\usepackage[scr=rsfs,cal=euler]{mathalfa}
\usepackage{xfrac}
\usepackage[unicode,hidelinks,bookmarks=false]{hyperref}
\newcommand{\ie}{i.e.\ }
\newcommand{\cf}{cf.\ }
\newcommand{\resp}{respectively\ }
\newcommand{\Subsection}{\subsection}
\providecommand{\nobreakdash}{\nobreak\hbox{-}}
\setlength{\emergencystretch}{2em}
\multlinegap0pt
\usepackage{amssymb}
\usepackage{mathtools}
\usepackage[shortlabels]{enumitem}
\newtheorem{proposition}{Proposition}
\title{de la Vall\'ee Poussin Means of Walsh-Fourier Expansions}
\author{Ushangi Goginava}
\date{Source: arXiv:2605.05135v1 (CC BY 4.0)}
\begin{document}
\maketitle

\noindent\emph{Source and licence.} de la Vall\'ee Poussin Means of Walsh-Fourier Expansions. Version arXiv:2605.05135v1 (CC BY 4.0), distributed by arXiv under the Creative Commons Attribution 4.0 International licence, http://creativecommons.org/licenses/by/4.0/, as recorded in the arXiv licence metadata for this exact version and shown on the abstract page https://arxiv.org/abs/2605.05135v1. Full licence notice: \texttt{licenses/arxiv-2605.05135-CC-BY-4.0.txt}.\par\smallskip

\noindent\textit{Editorial scope.} Counted unit: the article's proposition prop:block (source0.tex lines 199--234). The article's complete original proof of it is delivered with it (source0.tex lines 236--369, one \texttt{proof} environment of 134 lines). No mathematics is written, repaired or re-derived: the statement and the proof are the article's own lines, in the article's own notation, and no retained line is substituted or re-flowed.  No further source-local context is needed: every cross-reference of every retained line resolves inside the two delivered spans.  Nothing was omitted from any retained span, so the package carries no omission record.  No retained line contains a citation, so the wrapper carries no bibliography and no bibliography entry.  No retained line contains a figure environment, an image-inclusion command or drawing code, so no image is delivered and the package declares no asset dependency.  Editorial formatting only, in addition: an AMS \texttt{amsart} preamble that restates the article's own declarations and macros used by the retained lines, namely the environments \texttt{proposition} on separate counters, together with the standard packages those lines need.  The retained environments print in this document as Proposition 1 in order of appearance, whereas the article prints the same items with the numbering of its own full document.  The complete native source of the article as distributed in the arXiv e-print for this exact version is bundled unchanged as \texttt{sources/199926/source0.tex}.
\begin{proposition}
\label{prop:block} Let $m,\gamma\in\mathbb{N}$ with $2\gamma<m$. Define 
\begin{equation}  \label{eq:def-mu}
\mu(v,j):=2^{m-\gamma}v+2^{m-2\gamma}(v\oplus 2^j), \qquad 0\le
v<2^{\gamma}, \quad 0\le j<\gamma,
\end{equation}
and 
\begin{equation}  \label{eq:def-P}
P_{m,\gamma}(x):=\frac{1}{\sqrt{\gamma}}\sum_{v=0}^{2^{\gamma}-1}%
\sum_{j=0}^{\gamma-1} w_{\mu(v,j)}(x).
\end{equation}
Then the following assertions hold.

\begin{itemize}
\item[(i)] $\|P_{m,\gamma}\|_{L^{1}([0,1))}\le1$.

\item[(ii)] $\|P_{m,\gamma}\|_{L^{\infty}([0,1))}\le 2^{\gamma}\sqrt{\gamma}$%
.

\item[(iii)] 
\begin{equation*}
\mathrm{Spec}(P_{m,\gamma})\subset [2^{m-2\gamma},2^m)\cap 2^{m-2\gamma}%
\mathbb{N}_0.
\end{equation*}

\item[(iv)] For every $x\in[0,1)$ there exists an integer $\ell(x)$ such
that 
\begin{equation*}
\ell(x)\in [2^{m-\gamma},2^m)\cap 2^{m-2\gamma}\mathbb{N}
\end{equation*}
and 
\begin{equation*}
|S_{\ell(x)}(P_{m,\gamma};x)|\ge \frac14\sqrt{\gamma}.
\end{equation*}
\end{itemize}
\end{proposition}

\begin{proof}
Let 
\begin{equation*}
v=\sum_{k=0}^{\gamma -1}\nu _{k}2^{k},\qquad \nu _{k}\in \{0,1\}.
\end{equation*}%
Since the binary digit blocks of $2^{m-\gamma }v$ and $2^{m-2\gamma
}(v\oplus 2^{j})$ are disjoint, we have 
\begin{equation*}
w_{\mu (v,j)}(x)=w_{2^{m-\gamma }v}(x)\,w_{2^{m-2\gamma }(v\oplus 2^{j})}(x).
\end{equation*}%
A direct computation gives 
\begin{equation*}
w_{2^{m-\gamma }v}(x)=\prod_{k=0}^{\gamma -1}r_{m-\gamma +k}(x)^{\nu _{k}}
\end{equation*}%
and 
\begin{equation*}
w_{2^{m-2\gamma }(v\oplus 2^{j})}(x)=r_{m-2\gamma +j}(x)\prod_{k=0}^{\gamma
-1}r_{m-2\gamma +k}(x)^{\nu _{k}}.
\end{equation*}%
Therefore 
\begin{eqnarray*}
w_{\mu (v,j)}(x) &=&r_{m-2\gamma +j}(x)\prod_{k=0}^{\gamma -1}\bigl(%
r_{m-\gamma +k}(x)r_{m-2\gamma +k}(x)\bigr)^{\nu _{k}} \\
&=&r_{m-2\gamma +j}(x)c\left( x\right) ,
\end{eqnarray*}%
where%
\begin{equation*}
c\left( x\right) :=\prod_{k=0}^{\gamma -1}\bigl(r_{m-\gamma
+k}(x)r_{m-2\gamma +k}(x)\bigr)^{\nu _{k}}\in \left\{ -1,1\right\} .
\end{equation*}%
Summing over $v$ yields 
\begin{equation}
P_{m,\gamma }(x)=\frac{1}{\sqrt{\gamma }}\sum_{j=0}^{\gamma -1}r_{m-2\gamma
+j}(x)\prod_{k=0}^{\gamma -1}\bigl(1+r_{m-\gamma +k}(x)r_{m-2\gamma +k}(x)%
\bigr).  \label{eq:P-factor}
\end{equation}

Let 
\begin{equation*}
E_{m,\gamma}:=\bigl\{x\in[0,1): r_{m-\gamma+k}(x)=r_{m-2\gamma+k}(x)\text{
for }0\le k<\gamma\bigr\}.
\end{equation*}
Then \eqref{eq:P-factor} becomes 
\begin{equation}  \label{eq:P-on-E}
P_{m,\gamma}(x)=\frac{2^{\gamma}}{\sqrt{\gamma}}\mathbf{1}%
_{E_{m,\gamma}}(x)\sum_{j=0}^{\gamma-1} r_{m-2\gamma+j}(x).
\end{equation}
Assertion~(ii) follows immediately from \eqref{eq:P-on-E}.

To prove~(i), note that on $E_{m,\gamma}$ the vector 
\begin{equation*}
\bigl(r_{m-2\gamma}(x),\dots,r_{m-\gamma-1}(x)\bigr)
\end{equation*}
assumes each value in $\{\pm1\}^{\gamma}$ on a set of measure $2^{-2\gamma}$%
. Hence, by \eqref{eq:P-on-E}, 
\begin{align*}
\|P_{m,\gamma}\|_{L^{1}([0,1))} &=\frac{2^{\gamma}}{\sqrt{\gamma}}%
\int_{E_{m,\gamma}} \left|\sum_{j=0}^{\gamma-1} r_{m-2\gamma+j}(x)\right|\,dx
\\
&=\frac{1}{2^{\gamma}\sqrt{\gamma}}\sum_{\varepsilon\in\{\pm1\}^{\gamma}}
|\varepsilon_1+\cdots+\varepsilon_{\gamma}| \\
&\le \frac{1}{\sqrt{\gamma}}\left(\frac{1}{2^{\gamma}}\sum_{\varepsilon\in\{%
\pm1\}^{\gamma}} |\varepsilon_1+\cdots+\varepsilon_{\gamma}|^2\right)^{1/2}
\le 1.
\end{align*}
This proves~(i).

Since every $\mu(v,j)$ is a multiple of $2^{m-2\gamma}$, the right-hand side
of \eqref{eq:def-P} shows that 
\begin{equation*}
\mathrm{Spec}(P_{m,\gamma})\subset 2^{m-2\gamma}\mathbb{N}_0.
\end{equation*}
Moreover, 
\begin{equation*}
2^{m-2\gamma}\le \mu(v,j)<2^m,
\end{equation*}
which proves~(iii).

Fix $x\in[0,1)$. Choose $\sigma(x)\in\{\pm1\}$ that occurs at least $\gamma
/2$ times among 
\begin{equation*}
r_{m-2\gamma}(x),\dots,r_{m-\gamma-1}(x),
\end{equation*}
and put 
\begin{equation*}
J(x):=\{0\le j<\gamma: r_{m-2\gamma+j}(x)=\sigma(x)\}, \qquad \#J(x)\ge 
\frac{\gamma}{2},
\end{equation*}
\begin{equation*}
v(x):=\sum_{j\in J(x)} 2^j.
\end{equation*}
Define 
\begin{equation*}
\ell_1(x):=2^{m-\gamma}v(x), \qquad
\ell_2(x):=2^{m-\gamma}v(x)+2^{m-2\gamma}v(x).
\end{equation*}
Since $v(x)\ge1$, we have 
\begin{equation*}
2^{m-\gamma}\le \ell_1(x)<\ell_2(x)<2^m, \qquad \ell_1(x),\ell_2(x)\in
2^{m-2\gamma}\mathbb{N}.
\end{equation*}
For fixed $v=v(x)$ we have 
\begin{equation*}
\mu(v,j)=2^{m-\gamma}v+2^{m-2\gamma}(v\oplus 2^j).
\end{equation*}
Because $v\oplus 2^j<v$ if and only if the $j$th binary digit of $v$ equals $%
1$, it follows that 
\begin{equation*}
\mu(v(x),j)\in [\ell_1(x),\ell_2(x)) \qquad \Longleftrightarrow \qquad j\in
J(x).
\end{equation*}
Consequently, 
\begin{align*}
&S_{\ell_2(x)}(P_{m,\gamma};x)-S_{\ell_1(x)}(P_{m,\gamma};x) \\
&\qquad = \frac{1}{\sqrt{\gamma}}\sum_{j\in J(x)} w_{\mu(v(x),j)}(x) =\frac{%
c(x)}{\sqrt{\gamma}}\sum_{j\in J(x)} r_{m-2\gamma+j}(x)
\end{align*}
with a unimodular factor $c(x)\in\{\pm1\}$ independent of $j$. Since $%
r_{m-2\gamma+j}(x)=\sigma(x)$ for every $j\in J(x)$ and $|c(x)|=1$, it
follows that 
\begin{equation*}
\left|S_{\ell_2(x)}(P_{m,\gamma};x)-S_{\ell_1(x)}(P_{m,\gamma};x)\right|= 
\frac{\#J(x)}{\sqrt{\gamma}}\ge \frac12\sqrt{\gamma}.
\end{equation*}
Since for arbitrary real numbers $A$ and $B$ one has 
\begin{equation*}
\max\{|A|,|A+B|\}\ge \frac12|B|,
\end{equation*}
we conclude that at least one of $\ell_1(x)$ and $\ell_2(x)$ satisfies 
\begin{equation*}
|S_{\ell(x)}(P_{m,\gamma};x)|\ge \frac14\sqrt{\gamma}.
\end{equation*}
This proves~(iv).
\end{proof}

\end{document}
